% 分野別類題: 熱伝導方程式（変数分離法） 解答例
% コンパイル: TEXINPUTS="../../../00_common:$TEXINPUTS" latexmk -xelatex answers.tex
\documentclass[a4paper,11pt]{article}
\usepackage{preamble}
\usepackage{shoakubox}

\begin{document}

\kamokumei{類題演習 解答例 --- 熱伝導方程式（変数分離法）}

\daimon{【解法】熱伝導方程式の変数分離法による解き方と導出}

\noindent
有限区間 $0 < x < L$ における熱伝導方程式（熱方程式ともいう。同じ方程式の呼び名の違いである）
\[
u_t = k\,u_{xx},
\quad\text{すなわち}\quad
\frac{\partial u}{\partial t} = k\,\frac{\partial^{2} u}{\partial x^{2}}
\qquad (k > 0 \text{ は熱拡散率})
\]
を、境界条件のもとで変数分離法により解く一般的な手順を示す。ここで $u_t$, $u_{xx}$ は偏微分の略記であり、それぞれ $\dfrac{\partial u}{\partial t}$, $\dfrac{\partial^{2} u}{\partial x^{2}}$ を表す。

\medskip\noindent
\textbf{1. 変数分離}\quad $u(x,t) = X(x)\,T(t)$ とおいて方程式に代入すると
\[
X(x)\,T'(t) = k\,X''(x)\,T(t)
\quad\Longrightarrow\quad
\frac{T'(t)}{k\,T(t)} = \frac{X''(x)}{X(x)}
\]
左辺は $t$ のみ、右辺は $x$ のみの関数であるから、両辺は定数に等しい。これを $-\lambda$（$\lambda$ は分離定数）とおくと
\[
X''(x) + \lambda X(x) = 0, \qquad T'(t) + k\lambda\,T(t) = 0
\]

\medskip\noindent
\textbf{2. 固有値問題（境界条件による $\lambda$ の決定）}

\noindent
\textit{(a) ディリクレ境界条件 $u(0,t) = u(L,t) = 0$ の場合：}\quad $X(0) = X(L) = 0$ が要求される。$\lambda \le 0$ では自明解しか得られないことが確かめられるので $\lambda = \mu^2 > 0$ とおくと、$X(x) = A\cos\mu x + B\sin\mu x$ のうち $X(0)=0$ より $A=0$、さらに $X(L) = B\sin\mu L = 0$（$B \neq 0$）より
\[
\sin\mu L = 0 \quad\Longrightarrow\quad \mu_n = \frac{n\pi}{L}, \qquad \lambda_n = \left(\frac{n\pi}{L}\right)^{2} \qquad (n = 1, 2, 3, \dots)
\]
すなわち固有関数は $X_n(x) = \sin\dfrac{n\pi x}{L}$ である。

\noindent
\textit{(b) ノイマン境界条件 $u_x(0,t) = u_x(L,t) = 0$ の場合：}\quad $X'(0) = X'(L) = 0$ が要求される。同様に $X(x) = A\cos\mu x + B\sin\mu x$ とすると $X'(x) = -A\mu\sin\mu x + B\mu\cos\mu x$ より $X'(0) = B\mu = 0$ で $B = 0$（$\mu \neq 0$ のとき）、$X'(L) = -A\mu\sin\mu L = 0$（$A \neq 0$）より
\[
\sin\mu L = 0 \quad\Longrightarrow\quad \mu_n = \frac{n\pi}{L}, \qquad \lambda_n = \left(\frac{n\pi}{L}\right)^{2} \qquad (n = 0, 1, 2, \dots)
\]
固有関数は $X_n(x) = \cos\dfrac{n\pi x}{L}$（$n=0$ のとき定数 $X_0 = 1$ も許される）。

\medskip\noindent
\textbf{3. 時間成分}\quad $T'(t) + k\lambda_n T(t) = 0$ を解いて $T_n(t) = e^{-k\lambda_n t}$。

\medskip\noindent
\textbf{4. 重ね合わせと初期条件}\quad 各 $u_n(x,t) = X_n(x)\,T_n(t)$ は方程式と境界条件を満たすので、一般解は
\[
u(x,t) = \sum_{n} c_n\,X_n(x)\,e^{-k\lambda_n t}
\]
初期条件 $u(x,0) = f(x)$ を課すと $f(x) = \sum_n c_n X_n(x)$ となり、これは $f$ のフーリエ正弦級数（ディリクレの場合）またはフーリエ余弦級数（ノイマンの場合）に他ならない。係数はそれぞれ
\[
c_n = \frac{2}{L}\int_0^L f(x)\sin\frac{n\pi x}{L}\,dx \quad(\text{ディリクレ}), \qquad
c_n = \frac{2}{L}\int_0^L f(x)\cos\frac{n\pi x}{L}\,dx\ (n\ge1),\ \ c_0 = \frac{1}{L}\int_0^L f(x)\,dx
\]
で与えられる。$f(x)$ がすでに固有関数の有限和で書けている場合は、係数比較のみで $c_n$ が直ちに定まる（フーリエ係数の計算を要しない）。

\clearpage
\begin{mondaiboxS}{問1}
\[
u_t = k\,u_{xx}\ (0<x<L,\ t>0), \quad u(0,t)=u(L,t)=0, \quad u(x,0) = \sin\frac{\pi x}{L}
\]
を解け。
\end{mondaiboxS}
\kaitou
\textbf{手順1: 変数分離。}\ $u(x,t) = X(x)\,T(t)$ とおいて $u_t = k\,u_{xx}$ に代入すると
\[
X(x)\,T'(t) = k\,X''(x)\,T(t)
\]
両辺を $k\,X(x)\,T(t)$ で割ると
\[
\frac{T'(t)}{k\,T(t)} = \frac{X''(x)}{X(x)}
\]
左辺は $t$ のみ、右辺は $x$ のみの関数なのに両者が等しいので、その値は $x$ にも $t$ にも依存しない定数である。これを $-\lambda$ とおくと、2本の常微分方程式
\[
X''(x) + \lambda X(x) = 0, \qquad T'(t) + k\lambda\,T(t) = 0
\]
に分離される。

\textbf{手順2: 境界条件を $X$ に移す。}\ $u(0,t) = X(0)\,T(t) = 0$ がすべての $t$ で成り立つには、$T \equiv 0$（自明解）でない限り $X(0) = 0$。同様に $X(L) = 0$。

\textbf{手順3: 固有値問題を解く（$\lambda$ の場合分け）。}

\textbf{(i) $\lambda < 0$（$\lambda = -\mu^{2}$, $\mu>0$）のとき。}\\
一般解は
\[
X = A\cosh(\mu x) + B\sinh(\mu x)
\]
境界条件を課すと
\begin{align*}
X(0) &= A = 0 \\
X(L) &= B\sinh(\mu L) = 0
\end{align*}
$\sinh(\mu L) \neq 0$ より $B = 0$。自明解のみ。

\textbf{(ii) $\lambda = 0$ のとき。}\\
一般解は $X = A + Bx$。$X(0) = A = 0$、$X(L) = BL = 0$ より $B = 0$。自明解のみ。

\textbf{(iii) $\lambda > 0$（$\lambda = \mu^{2}$, $\mu>0$）のとき。}\\
一般解は
\[
X = A\cos(\mu x) + B\sin(\mu x)
\]
境界条件を課すと
\begin{align*}
X(0) &= A = 0 \\
X(L) &= B\sin(\mu L) = 0
\end{align*}
自明でない解（$B \neq 0$）が存在するためには $\sin(\mu L) = 0$、すなわち
\[
\mu_n = \frac{n\pi}{L}, \qquad \lambda_n = \left(\frac{n\pi}{L}\right)^{2}, \qquad X_n(x) = \sin\frac{n\pi x}{L} \qquad (n = 1, 2, 3, \dots)
\]

\textbf{手順4: 時間成分を解く。}\ $T' = -k\lambda_n T$ は変数分離形で
\[
\frac{dT}{T} = -k\lambda_n\,dt
\quad\Longrightarrow\quad
\ln|T| = -k\lambda_n t + C
\quad\Longrightarrow\quad
T_n(t) = e^{-k\lambda_n t}
\]
（定数倍は後の $c_n$ に吸収させる。）

\textbf{手順5: 重ね合わせと初期条件。}\ 各 $X_n(x)T_n(t)$ は方程式と境界条件を満たすので、線形性からその重ね合わせ
\[
u(x,t) = \sum_{n=1}^{\infty} c_n \sin\frac{n\pi x}{L}\,e^{-k(n\pi/L)^2 t}
\]
も満たす。$t=0$ とおくと
\[
u(x,0) = \sum_{n=1}^{\infty} c_n \sin\frac{n\pi x}{L} = \sin\frac{\pi x}{L}
\]
右辺はすでに $n=1$ の固有関数そのものなので、係数比較により $c_1 = 1$、$c_n = 0\ (n \ge 2)$（固有関数系 $\{\sin\frac{n\pi x}{L}\}$ の直交性から、この対応は一意である）。よって
\[
\boxed{\; u(x,t) = \sin\frac{\pi x}{L}\,\exp\!\left(-k\left(\frac{\pi}{L}\right)^{2} t\right) \;}
\]
（検算: $u_t = -k(\pi/L)^2 \sin\frac{\pi x}{L}e^{-k(\pi/L)^2t}$、$u_{xx} = -(\pi/L)^2\sin\frac{\pi x}{L}e^{-k(\pi/L)^2t}$ より $u_t = k u_{xx}$ が成り立ち、$u(0,t)=u(L,t)=0$、$u(x,0)=\sin\frac{\pi x}{L}$ も満たす。）

\clearpage
\begin{mondaiboxS}{問2}
\[
u_t = k\,u_{xx}\ (0<x<L,\ t>0), \quad u(0,t)=u(L,t)=0, \quad u(x,0) = 3\sin\frac{\pi x}{L} - 2\sin\frac{3\pi x}{L}
\]
を解け。
\end{mondaiboxS}
\kaitou
境界条件が問1と同じディリクレ型なので、変数分離から固有値問題までの手順（問1の手順1〜4）はそのまま使え、一般解は
\[
u(x,t) = \sum_{n=1}^\infty c_n \sin\frac{n\pi x}{L}\,e^{-k(n\pi/L)^2 t}
\]
である。$t = 0$ とおいて初期条件と比べると
\[
\sum_{n=1}^\infty c_n \sin\frac{n\pi x}{L} = 3\sin\frac{\pi x}{L} - 2\sin\frac{3\pi x}{L}
\]
右辺はすでに固有関数 $\sin\dfrac{\pi x}{L}$（$n=1$）と $\sin\dfrac{3\pi x}{L}$（$n=3$）の線形結合なので、フーリエ係数の積分計算をするまでもなく係数比較で
\[
c_1 = 3, \qquad c_3 = -2, \qquad c_n = 0\ (n \neq 1, 3)
\]
と定まる（固有関数系の直交性により、この表し方は一意）。各モードは自分の固有値に応じた速さ $e^{-k(n\pi/L)^2 t}$ で独立に減衰するから、それぞれの係数に指数因子を付ければよい。よって
\[
\boxed{\; u(x,t) = 3\sin\frac{\pi x}{L}\,\exp\!\left(-k\left(\frac{\pi}{L}\right)^{2} t\right) - 2\sin\frac{3\pi x}{L}\,\exp\!\left(-k\left(\frac{3\pi}{L}\right)^{2} t\right) \;}
\]
（検算: 各項は問1と同じ形の固有モードであり、線形性から $u_t = k u_{xx}$ を満たす。境界: $x=0,L$ で $\sin(n\pi x/L)=0$ なので $u(0,t)=u(L,t)=0$。初期条件: $t=0$ で $e^{0}=1$ となり $u(x,0)=3\sin\frac{\pi x}{L}-2\sin\frac{3\pi x}{L}$ に一致。）

\clearpage
\begin{mondaiboxS}{問3}
\[
u_t = k\,u_{xx}\ (0<x<L,\ t>0), \quad u_x(0,t)=u_x(L,t)=0, \quad u(x,0) = u_0 + \cos\frac{\pi x}{L}
\]
（両端断熱、$u_0$ は定数）を解け。
\end{mondaiboxS}
\kaitou
\textbf{手順1: 変数分離。}\ 問1と同じく $u = X(x)T(t)$ とおくと
\[
X'' + \lambda X = 0, \qquad T' + k\lambda T = 0
\]
に分離される。今回はノイマン境界条件なので、$u_x(0,t) = X'(0)T(t) = 0$ より $X'(0) = 0$、同様に $X'(L) = 0$ が $X$ に課される。

\textbf{手順2: 固有値問題を解く（$\lambda$ の場合分け）。}

\textbf{(i) $\lambda < 0$（$\lambda = -\mu^{2}$, $\mu>0$）のとき。}\\
一般解とその導関数は
\[
X = A\cosh(\mu x) + B\sinh(\mu x), \qquad
X' = A\mu\sinh(\mu x) + B\mu\cosh(\mu x)
\]
境界条件を課すと
\begin{align*}
X'(0) &= B\mu = 0 \\
X'(L) &= A\mu\sinh(\mu L) = 0
\end{align*}
$\mu > 0$ より $B = 0$、$\sinh(\mu L) \neq 0$ より $A = 0$。自明解のみ。

\textbf{(ii) $\lambda = 0$ のとき。}\\
一般解とその導関数は $X = A + Bx$、$X' = B$。$X'(0) = X'(L) = B = 0$ とすれば $A$ は任意でよく、\textbf{定数解 $X_0 = 1$ が自明でない解として生き残る}（ディリクレの場合との最大の違い）。$\lambda_0 = 0$ は固有値である。

\textbf{(iii) $\lambda > 0$（$\lambda = \mu^{2}$, $\mu>0$）のとき。}\\
一般解とその導関数は
\[
X = A\cos(\mu x) + B\sin(\mu x), \qquad
X' = -A\mu\sin(\mu x) + B\mu\cos(\mu x)
\]
境界条件を課すと
\begin{align*}
X'(0) &= B\mu = 0 \\
X'(L) &= -A\mu\sin(\mu L) = 0
\end{align*}
$B = 0$、そして自明でない解（$A \neq 0$）のためには $\sin(\mu L) = 0$、すなわち
\[
\mu_n = \frac{n\pi}{L}, \qquad \lambda_n = \left(\frac{n\pi}{L}\right)^{2}, \qquad X_n(x) = \cos\frac{n\pi x}{L} \qquad (n = 1, 2, 3, \dots)
\]

\textbf{手順3: 時間成分と重ね合わせ。}\ $T_n(t) = e^{-k\lambda_n t}$（問1の手順4と同じ）。$n=0$ のモードは $\lambda_0 = 0$ なので $T_0(t) = e^{0} = 1$、つまり\textbf{減衰せずに残る}。重ね合わせて一般解は
\[
u(x,t) = c_0 + \sum_{n=1}^{\infty} c_n \cos\frac{n\pi x}{L}\,e^{-k(n\pi/L)^2 t}
\]

\textbf{手順4: 初期条件。}\ $t=0$ とおくと
\[
c_0 + \sum_{n=1}^{\infty} c_n \cos\frac{n\pi x}{L} = u_0 + \cos\frac{\pi x}{L}
\]
右辺はすでに $n=0$ の定数モードと $n=1$ の固有関数の和なので、係数比較により $c_0 = u_0$、$c_1 = 1$、他の $c_n = 0$。よって
\[
\boxed{\; u(x,t) = u_0 + \cos\frac{\pi x}{L}\,\exp\!\left(-k\left(\frac{\pi}{L}\right)^{2} t\right) \;}
\]
（検算: $u_t = -k(\pi/L)^2\cos\frac{\pi x}{L}e^{-k(\pi/L)^2t}$、$u_{xx} = -(\pi/L)^2\cos\frac{\pi x}{L}e^{-k(\pi/L)^2t}$ より $u_t = k u_{xx}$。$u_x(x,t) = -\frac{\pi}{L}\sin\frac{\pi x}{L}e^{-k(\pi/L)^2t}$ は $x=0,L$ で $\sin 0 = \sin\pi = 0$ となり $u_x(0,t)=u_x(L,t)=0$ を満たす。初期条件も $t=0$ で一致。また $t\to\infty$ で $u \to u_0$ となり、断熱系ではエネルギー（区間平均温度）が保存され定常状態が初期の平均値 $u_0$ に一致することとも整合する。）

\end{document}
